\documentclass[
	xcolor=dvipsnames,
	10pt, 
	]{beamer}
\graphicspath{{./fig/}}
\RequirePackage{"./sty/UserPackageManager"}
\usepackage{changepage}
%% TUBS
  \usebackgroundtemplate%
{
	\includegraphics[width=\paperwidth,height=\paperheight]{./bonos/Bonobono}%
}

\usepackage{biblatex}
\bibliography{Bibilography.bib}
\title[]{Notes on Kinematic Fit}
\author[Author]{Bono Bono}
\institute{Bono University}
\date{\today}
\subtitle{Bono, Bono\footnotemark[1]\newline\today}


\begin{document}

%%%%%%%%%%%%%%% Title page 
\begin{frame}[t,plain] % Cover slide
       \titlepage
        \footnotetext[1]{Bono University}
\end{frame}
\begin{frame}{Introduction}
	\begin{block}{}
	Assume that we have measured a momentum of two particles, which decay from a mother particle with an exact momentum $P_0$. In the real world, every measurement inherently carries some errors. Consequently, the measured momentum may not satisfy momentum conservation:
	\begin{align}
	P_{0}=P_{1}+P_{2};\quad P_{1,meas}+P_{2,meas}\neq P_{0} 
	\end{align}
	However, since these measurements did not incorporate our prior knowledge from physics, we can make a more informed estimate of the measured parameters. If the momentum resolution of particles 1 and 2 is well-known, then we can express $\chi^2$ as
	\begin{align}
		\chi^2= \frac{(P_{1}-P_{1,meas})^2}{\sigma_1^2}+\frac{(P_{2}-P_{2,meas})^2}{\sigma_2^2}\label{chi2}
	\end{align}
	\end{block}
\end{frame}
\begin{frame}{Introduction}
	\begin{block}{}
		By incorporating the \textit{Kinematic Constraints}, specifically \textit{momentum conservation}, into our example, we introduce additional terms known as \textit{Lagrange Multiplier} to Equation \eqref{chi2}:
		\vspace{-5 mm}
		\begin{align}
			\chi^2= \frac{(P_{1}-P_{1,meas})^2}{\sigma_1^2}+\frac{(P_{2}-P_{2,meas})^2}{\sigma_2^2}\label{chi2} + 2\mathbf{\lambda(P_{1,meas}+P_{2,meas}-P_0)}
		\end{align} 
		We then proceed to evaluate the conditions for local minima, i.e. setting the partial derivatives equal to zero:
		\begin{align}
			\frac{1}{2}\pdf{\chi2}{P_{1}}&=\frac{(P_{1}-P_{1,meas})}{\sigma_1^2}+\lambda=0\label{dxd1}\\
			\frac{1}{2}\pdf{\chi2}{P_{2}}&=\frac{(P_{2}-P_{2,meas})}{\sigma_2^2}+\lambda=0\label{dxd2}
			\\
			\frac{1}{2}\pdf{\chi2}{\lambda}&=(P_{1}+P_{2}-P_0)=0\label{dxdl}
		\end{align}
	\end{block}
\end{frame}
\begin{frame}{Introduction}
	\begin{block}{}
		By solving the equations \ref{dxd1},\ref{dxd2},\ref{dxdl}, we obtain the following expressions:
		\begin{align}
			P_1& =\frac{\sigma_2^2 P_{1,meas}-\sigma_1^2P_{2,meas}+\sigma_1^2P_0}{\sigma_1^2+\sigma_2^2}\\
			P_2& = \frac{\sigma_1^2 P_{2,meas}-\sigma_2^2P_{1,meas}+\sigma_2^2P_0}{\sigma_1^2+\sigma_2^2}\\
			\lambda&=\frac{P_{1,meas}+P_{2,meas}-P_0}{\sigma_1^2+\sigma_2^2}.
		\end{align}
			Now, we have obtained the 'corrected' measurements with minimized $\chi^2$, which incorporates momentum conservation.
			Let us delve into the interpretation of these equations.
	\end{block}
\end{frame}
\begin{frame}{Introduction}
	\begin{block}{}
			 In a straightforward interpretation, $\lambda$ can be viewed as a kind of 'normalized variance' of the kinematic constraint. It quantifies the error in momentum conservation($P_{1,meas}+P_{2,meas}-P_0$) relative to the overall resolution($\sigma_1^2+\sigma_2^2$).  Equation \eqref{dxd1} implies that 
\begin{align}
	P_{1}=P_{1,meas}-\sigma_1^2\lambda,
\end{align}
			suggesting that the corrected momentum($P_1$) is essentially the measured momentum($P_{1,meas}$) augmented by a term proportional to the detector resolution and the normalized error of the kinematic constraint. Thus, we can assert that we have applied a statistically fair correction to the momentum, taking into account both the detector resolution and kinematic constraints.
	\end{block}
\end{frame}
\begin{frame}{Fitting in General}
	\begin{block}{}
		Assume that you have a set of measurements, $\textbf m =\{m_1,m_2\ldots m_N\}$, and some unmeasured data, $\textbf u =\{u_1,u_2\ldots u_J\}$ to be estimated. Kinematic constraints can be represented by sets of equations $\textbf f =\{f_1(m_1,m_2,\ldots m_N,u_1,u_2,\ldots u_N),f_2,\ldots f_K\}$. We will iteratively solve the problem by guessing the best parameter for each step and checking  $\chi^2$.  Let $\mathbf m^0$ denote our initial measured data, and $\mathbf m$ represent the 'guess' of the data in each iterative step.
		\vspace{-2 mm}
		\begin{align}
			\chi^2(\textbf{m}) = (\textbf{m}^0-\textbf{m})^\dagger V^{-1}(\textbf{m}^0-\textbf{m})+2\mathbf\lambda^\dagger \textbf f(\textbf{m,u}).\label{KFChi2}
		\end{align}
		Here, the Lagrange multiplier $\mathbf \lambda =\{\lambda_1,\lambda_2,\ldots\lambda_K\}$ is not just a number but a column vector with k elements, corresponding to each kinematic constraint in $\mathbf f$. Our task is to minimize $\chi^2$ to obtain the best guesses in statistically fair method.	
	\end{block}
\end{frame}
\begin{frame}{$\chi^2$ Minimization}
	\begin{block}{}
		By (partially)differentiating with respect to all variables involved, we obtain the gradients of $\chi^2$. Setting all of them to zero indicates that we have reached a minimum point of $\chi^2$. We have 3 sets of gradient equations:
		\begin{align}
			\nabla_{\mathbf m} &= -2 V^{-1}(\mathbf m^0)(\mathbf m^0-\mathbf m) + 2 \mathbf F_\mathbf m^\dagger(\mathbf{m,u})\mathbf\lambda=0\label{GradM}\\
			\nabla_{\mathbf u} &= 2\mathbf F_\mathbf u^\dagger(\mathbf{m,u})\mathbf \lambda =0\label{GradU}\\
			\nabla_{\mathbf \lambda}& = \mathbf f(\mathbf{m,u})\label{GradL}.
		\end{align}
		Here, the subscripts denote partial derivatives. i.e. ($(\mathbf {F}_m)_{ki}\equiv\pdf {f_k}{m_i}$).
	\end{block}
\end{frame}
\begin{frame}{Processing Iterative Steps.}
	\begin{block}{}
		We can express the following equations based on the ones provided above:
		\begin{align}
			&V^{-1}(\mathbf m^0)(\mathbf{m}^{\nu+1}-\mathbf{m}^0)+(\mathbf{F}_m^\dagger)^\nu\mathbf\lambda^{\nu+1}=0\label{DelM}\\
			&(\mathbf{F}_u^\dagger)^\nu\mathbf\lambda^{\nu+1}=0\label{DelU}\\
			&\mathbf f^\nu + \mathbf F^\nu_m (\mathbf m^{\nu+1}-\mathbf m^\nu) + \mathbf F^\nu_u(\mathbf u^{\nu+1}-\mathbf u^\nu)=0\label{DelL}.
		\end{align}
		Equation \eqref{DelL} is not a direct consequence of Equation \eqref{GradU} but rather a linear approximation to proceed with our iteration step. Expanding the $\nabla_\lambda$ term with a Taylor series leads to this equation. 
		Note that, as our parameters $\mathbf m$ and $\mathbf u$ are updated during the step, our constraint matrix $\mathbf f$ should also be updated during the iteration. Here, $\lambda$ should be indexed as $\nu+1$ since it is a parameter to be guessed in the next step.
	\end{block}
\end{frame}

\begin{frame}{Solving the Equation(1)}
	\begin{block}{}
		Multiplying \textbf V to Equation \eqref{DelM} leads to:		\vspace{-2 mm}
		\begin{align}
			\mathbf{m}^{\nu+1}-\mathbf{m}^0=-V(\mathbf m^0)(\mathbf{F}_m^\dagger)^\nu\mathbf\lambda^{\nu+1}.\label{Mnu}
		\end{align}
		Substituting Equation \eqref{Mnu} into Equation \eqref{DelL}, we get:\vspace{-2 mm}
		\begin{align}
			\mathbf F^\nu_u(\mathbf u^{\nu+1}-\mathbf u^\nu) &= - \mathbf f^\nu - \mathbf F^\nu_m (-V(\mathbf m^0)(\mathbf{F}_m^\dagger)^\nu\mathbf\lambda^{\nu+1} +\mathbf{m}^0-\mathbf{m}^\nu  )\nonumber\\
			&=  S\lambda^{\nu+1}  - R\label{SLambda}
		\end{align}
		where$S\equiv \mathbf{F}_\mathbf{m}^\nu V(\mathbf m^0)(\mathbf{F}_\mathbf{m}^\dagger )^\nu$  and $R\equiv \mathbf f^\nu + \mathbf{F}_\mathbf{m}^\nu(\mathbf{m}^0-\mathbf{m}^\nu)$. Multiplying $(\mathbf{F}_u^\dagger)^\nu S^{-1}$ and substituting Equation \eqref{DelU}, we get:\vspace{-2 mm}
		\begin{align}
			(\mathbf{F}_u^\dagger)^\nu S^{-1}\mathbf F^\nu_u(\mathbf u^{\nu+1}-\mathbf u^\nu) = \ctz{(\mathbf F_u^\dagger)^\nu\mathbf \lambda^{\nu+1}}-(\mathbf{F}_u^\dagger)^\nu S^{-1}R.
		\end{align}
	\end{block}
\end{frame}
\begin{frame}{Solving the Equation(2)}
	\begin{block}{}
		Then we naturally obtain:
		\begin{align}
			\mathbf{u}^{\nu+1}=\mathbf{u}^\nu - ((\mathbf{F}_u^\dagger)^\nu S^{-1}\mathbf F^\nu_u)^{-1}(\mathbf{F}_u^\dagger)^\nu S^{-1}R\label{Unu}.
		\end{align} 
		and from Equation \eqref{SLambda}
		\begin{align}
			\mathbf\lambda^{\nu+1} = S^{-1}(\mathbf F^\nu_u(\mathbf u^{\nu+1}-\mathbf u^\nu) + R)\label{Lnu}
		\end{align}.
		For a summary, we have obtained all equations to proceed to the next step. All other matrices in the equation can be calculated from parameters of the current step, and $\chi^2$ can be evaluated from \eqref{KFChi2} .
		\begin{align*}
			\begin{cases}
				\mathbf{u}^{\nu+1}=\mathbf{u}^\nu - ((\mathbf{F}_u^\dagger)^\nu S^{-1}\mathbf F^\nu_u)^{-1}(\mathbf{F}_u^\dagger)^\nu S^{-1}R&\eqref{Unu}\\
				\mathbf\lambda^{\nu+1} = S^{-1}(\mathbf F^\nu_u(\mathbf u^{\nu+1}-\mathbf u^\nu) + R)&\eqref{Lnu}\\
				\mathbf{m}^{\nu+1}=\mathbf{m}^0-V(\mathbf m^0)(\mathbf{F}_m^\dagger)^\nu\mathbf\lambda^{\nu+1}&\eqref{Mnu}
			\end{cases}
		\end{align*}
	\end{block}
\end{frame}
\begin{frame}{Evolution of Variance Matrix}
	\begin{block}{}
		Take a look at Eq.\eqref{Mnu}. We see that $m^{\nu+1} $ is an addition(subtraction) of some parameters to the initially measured data. As we already know the error, i.e. Variance matrix, of initial data, we can estimate how error propagates through the fitting process. 
		\begin{align}
			V(m) = J_{m,m^0}V(m^0)J_{m,m^0}^\dagger
		\end{align}
		We need to calculate the Jacobian, 
		\begin{align}
			J_{m,m^0(i,j)} = \pdf{m_i}{m^0_j}
		\end{align}
	\end{block}
\end{frame}
\begin{frame}{Evolution of Variance Matrix.}
	\begin{block}{}
		To begin with, let us express Eq \eqref{Mnu} in terms of $m^0$. At the moment we will drop the superscript $\nu$. As $\vecb f(\vecb m,\vecb u)$is a constant on $m^0$, $\vecb F_\vecb m$ also will be a constant to $m^0$. Then we only need to consider the derivatives of $\lambda$. By substituting \eqref{Unu} ,
		\begin{align}
	\mathbf\lambda = S^{-1}(-\mathbf F_u(((\mathbf{F}_u^\dagger) S^{-1}\mathbf F_u)^{-1}(\mathbf{F}_u^\dagger) S^{-1}R) + R)
		\end{align}
		and we have
		\begin{align}
			R\equiv \mathbf f + \mathbf{F}_\mathbf{m}(\mathbf{m}^0-\mathbf{m})\to \pdf{R}{m^0}= \vecb F_m
		\end{align} 
		so that 
		\begin{align}
			\pdf{\lambda}{m^0}= S^{-1}(-\mathbf F_u((\mathbf{F}_u^\dagger S^{-1}\mathbf F_u)^{-1}\mathbf{F}_u^\dagger S^{-1}\vecb F_m) + \vecb F_m).
		\end{align}
	\end{block}
\end{frame}
\begin{frame}{Evolution of Variance Matrix}
	\begin{block}{}
		Now define the symmetric matrices $G\equiv \vecb {F_m^\dagger} S^{-1}\vecb {F_m}$, $U\equiv (\mathbf{F}_u^\dagger S^{-1}\mathbf F_u)^{-1}$ and $H\equiv \mathbf{F}_m^\dagger S^{-1}\vecb {F_u}$. Then we have expressions for $\pdf{\lambda}{m^0}$ hence 
		\begin{align}
			J_{m,m^0}&=I - V(m^0)\vecb{F}_m^\dagger\pdf\lambda{m^0}=I - V\vecb{F}^\dagger_m(-S^{-1}\vecb F_u U^{-1}H^\dagger+S^{-1}\vecb F_m)\nonumber\\
			&=I-V(G-HUH^\dagger)
		\end{align}
		If we let $C =G-HUH^\dagger$, we obtain
		\begin{align}
			V(m) = J_{m,m^0} VJ_{m,m^0}^\dagger =V -2 VCV  +VCVCV.\label{Vm}
		\end{align}
		You might want to neglect higher order term, but please keep 2nd order term at the moment. Some materials like \cite{Prob} had an error in this part.
	\end{block}
\end{frame}
\begin{frame}{Variance of the Unknowns}
	\begin{block}{}
			Just like how we derived  Eq.\eqref{Vm} we can estimate the variance matrix of the unknowns. 
		\begin{align}
			V_U = J_{u,m0} V J_{u,m0}^T
		\end{align}
		$J_{u,m0}$ can be obtained from Eq.\eqref{Unu}. Denoting\\ $((\mathbf{F}_u^\dagger)^\nu S^{-1}\mathbf F^\nu_u)^{-1}(\mathbf{F}_u^\dagger)^\nu S^{-1}$ as $K$, 
		\begin{align}
			J_{u^{\nu+1},m^0} = \pdf {u^{\nu+1}} {m^0} = \pdf{u^{\nu}}{m^0} - K\pdf R {m^0} \simeq -K \vecb F_m.
		\end{align}
		We have initial "Guess" for the unknowns; In principle, it is not a driven value from measurements. Then, $\pdf {u^0}{m_0} = 0$. Also, we approximate that the terms in 2nd or higher iterations are negligible: $J_{u,m^0}\simeq J_{u^0,m^0}$.
	\end{block}
\end{frame}
\begin{frame}{Pull distribution}
	\begin{block}{}
		It is better to check the pull distribution to check the quality of a fit. By defining the residual $\epsilon = m-m^0$ and its variance $V(\epsilon)$, pull is defined as:, 
		\begin{align}
			P(\epsilon) = \epsilon/\sqrt{V(\epsilon)}
		\end{align}
		and
		\begin{align}
			V(\epsilon) \equiv V(m)+ V(m^0) - 2 Cov(m,m^0).\label{Ve}
		\end{align}
		We have already calculated $J_{m,m^0}$. Then we directly get the covariance matrix. 
	\end{block}
\end{frame} 
\begin{frame}{Pull distribution}
	\begin{block}{}
		The covariance can be estimated as:
		\begin{align}
			Cov(m,m^0) = J_{m,m^0}V(m) = V - VCV.
		\end{align}
		If we substitute this and Eq.\eqref{Vm} into Eq.\eqref{Ve}, we get
		\begin{align}
			V(\epsilon) = VCVCV.
		\end{align}
		Note that 2nd order term affects the variance.	
	\end{block}
\end{frame}
\begin{frame}{Example: 3-C Fit }
	\begin{block}{}
		Assuming a decay of $\Lambda\to p\pi^-$, let's represent the momentum of those particles in spherical coordinates:
		\begin{align}
			\vec{P}_\Lambda = (P_\Lambda,\theta_\Lambda,\phi_\lambda),\cdots
		\end{align}
		Then the kinematic constraints can be expressed as:
		\begin{align}
			\begin{pmatrix}
				f_1\\f_2\\f_3\\f_4
			\end{pmatrix}=
			\begin{pmatrix}
			-P_\Lambda\sin\theta_\Lambda\cos\phi_\Lambda+P_p\sin\theta_p\cos\phi_p+P_\pi\sin\theta_\pi\cos\phi_\pi\\
			-P_\Lambda\sin\theta_\Lambda\sin\phi_\Lambda+P_p\sin\theta_p\sin\phi_p+P_\pi\sin\theta_\pi\sin\phi_\pi\\
			-P_\Lambda\cos\theta_\Lambda+P_p\cos\theta_p+p_\pi\cos\theta_\pi\\
			-\sqrt{P_\Lambda^2+m_\Lambda^2}+\sqrt{P_p^2+m_p^2}+\sqrt{P_\pi^2+m_\pi^2}
			\end{pmatrix}.\label{FMat}
		\end{align}
	\end{block}
\end{frame}
\begin{frame}{Example: 3-C Fit}
	\begin{block}{}
		We have unmeasured and measured variables as:
		\begin{align}
			\mathbf{u} = \{P_\Lambda\};\quad \mathbf{m}=\{\theta_\Lambda,\phi_\Lambda,P_p,\theta_p,\phi_p,P_\pi,\theta_\pi,\phi_\pi\}.\label{variables}
		\end{align}		
		Since there is 1 unmeasured variable with 4 kinematical constraints, this is a 4-1 = 3-Constrained fit. Let us substitute Eq \eqref{FMat} and \eqref{variables} into Eq \eqref{KFChi2}s and its resulting equations. We get $\mathbf{F_u}$ and $\mathbf{F_m}$ as
		\begin{align}
			\mathbf{F_u}=
			\begin{pmatrix}
				\pdf{f_1}{P_\Lambda}\\
				\pdf{f_2}{P_\Lambda}\\
				\pdf{f_3}{P_\Lambda}\\
				\pdf{f_4}{P_\Lambda}
			\end{pmatrix};\quad \mathbf{F_m}=\begin{pmatrix}
			\pdf{f_1}{\theta_\Lambda} &\cdots&\pdf{f_1}{\phi_\pi} \\
			\pdf{f_2}{\theta_\Lambda} &\cdots&\pdf{f_2}{\phi_\pi}\\
			\pdf{f_3}{\theta_\Lambda} &\cdots&\pdf{f_3}{\phi_\pi}\\
			\pdf{f_4}{\theta_\Lambda} &\cdots&\pdf{f_4}{\phi_\pi}
			\end{pmatrix}
		\end{align}
		\begin{adjustwidth}{13mm}{}{
		We have all the matrices to calculate in each step. By applying an appropriate variance matrix and employing $\chi^2$ selection criteria, we can kinematically fit the particles.}
		\end{adjustwidth}
	\end{block}
\end{frame}
%\begin{frame}[allowframebreaks]{References}
\begin{frame}{References}
	 	\printbibliography
\end{frame}
\end{document}
